\documentclass[10pt]{article}

\usepackage[letterpaper,margin=1in]{geometry}
\usepackage{amsmath,amssymb,amsthm,mathtools}
\usepackage{booktabs}
\usepackage{enumitem}
\usepackage{microtype}
\usepackage{xcolor}
\usepackage[hidelinks]{hyperref}
\setlength{\emergencystretch}{1em}
\hypersetup{
  pdftitle={Solving the Shortest Vector Problem in half-exponential Time via Mid-point Hessian: Exact transport escapes and a response-floor obstruction},
  pdfauthor={Anonymous author}
}

\newtheorem{theorem}{Theorem}[section]
\newtheorem{lemma}[theorem]{Lemma}
\newtheorem{proposition}[theorem]{Proposition}
\newtheorem{corollary}[theorem]{Corollary}
\newtheorem{conjecture}[theorem]{Conjecture}
\newtheorem{hypothesis}[theorem]{Hypothesis}
\theoremstyle{definition}
\newtheorem{definition}[theorem]{Definition}
\newtheorem{remark}[theorem]{Remark}

\newcommand{\R}{\mathbb{R}}
\newcommand{\Z}{\mathbb{Z}}
\newcommand{\F}{\mathbb{F}}
\newcommand{\E}{\mathbb{E}}
\newcommand{\Prb}{\mathop{\rm Pr}}
\newcommand{\tr}{\mathop{\rm tr}}
\newcommand{\op}{\mathrm{op}}
\newcommand{\poly}{\mathrm{poly}}
\newcommand{\negl}{\mathrm{negl}}
\newcommand{\SVP}{\textnormal{\textsc{SVP}}}
\newcommand{\BDD}{\textnormal{\textsc{BDD}}}
\newcommand{\DGS}{\textnormal{\textsc{DGS}}}
\newcommand{\CVP}{\textnormal{\textsc{CVP}}}
\newcommand{\Haar}{\mathrm{Haar}}
\newcommand{\Torus}{\R^n/L}
\newcommand{\ip}[2]{\langle #1,#2\rangle}
\newcommand{\norm}[1]{\lVert #1\rVert}
\newcommand{\abs}[1]{\lvert #1\rvert}
\newcommand{\rhoo}{\rho}
\newcommand{\what}{\widehat}
\newcommand{\cT}{\mathcal{T}}
\newcommand{\cS}{\mathcal{S}}
\newcommand{\cG}{\mathcal{G}}
\newcommand{\cP}{\mathcal{P}}
\newcommand{\cE}{\mathcal{E}}

\title{Solving the Shortest Vector Problem in
$2^{0.5n+o(n)}$ Time\\via Mid-point Hessian\\
\large Exact transport escapes and a response-floor obstruction}
\author{Anonymous author\\[2mm]
\small Draft response to Hhan's mid-point Hessian algorithm}
\date{August 2026}

\begin{document}
\maketitle

\begin{abstract}
Hhan recently gave a randomized algorithm for exact shortest vector problem
running in time $2^{0.6039n+o(n)}$ and space $2^{0.5n+o(n)}$.  Its endpoint is
a geometric fact about the Hessian of the periodic Gaussian: at the midpoint
of a shortest vector, an extreme eigenvector points toward that shortest
vector and can be completed by preprocessing bounded distance decoding.

We study the information available before the Walsh--Hadamard search over the
$2^n$ parity midpoints.  A single panel of full-dual discrete Gaussian
samples gives the traceless periodic Hessian, and its first two derivatives,
at every point of the continuous torus.  We prove an exact half-grid
stationarity and snapping lemma, a bounded narrowing-importance identity,
and a uniform empirical derivative theorem using
$2^{(0.4002+o(1))n}$ samples at a rigorously controlled smooth scale.  A
log-trace-exponential objective removes every eigengap assumption from the
ascent vector field.  These results give a black-box reduction: any robust,
samplable ascent law reaching a shortest-vector parity with probability
$2^{-o(n)}$ yields a half-exponential algorithm.

We then resolve the proposed universal Haar-seeded law negatively.  For the
rectangular family
$L_n=\operatorname{diag}(1,2^n,\ldots,2^n)\Z^n$, the correct single scale is
$\Theta(n^{-1/2})$.  Outside exponentially small coordinate slabs, every
third derivative driving soft-Hessian ascent is
$\exp(-\Omega(n4^n))$.  Hence any polynomial-step schedule of polynomial bit
complexity leaves almost every Haar seed in its initial parity cell.  The
shortest parity is reached with probability at most $2^{-n+1}$, and the
robust basin is no larger.  This disproves the manuscript's former
Conjecture~13.1.  We then give exact lattice-adapted repairs.  Short dual
vectors impose syndrome-zero equations on every shortest parity; the
coefficient ellipsoid gives Hamming and spectral-tube covers; bounded
coefficient energy admits exact generalized-distributive-law elimination;
and an obtuse superbase reduces \(\SVP\) to one global minimum cut.  The last
clause shows that the all-ones cancellation fan is submodular, rather than a
genuine hard core.

For the surviving geometry we prove a causal transport formula.  An
infinitesimal isotropic displacement of a uniform point modulo the lattice
crosses a facet of parity \(\theta\) at rate proportional to that facet's
area.  The associated cone-volume weights form an exact isotropic frame, and
every shortest facet retains a constant-area robust core with explicit
inverse-polynomial reach and energy margin.  A separate gap theorem gives a
direct spherical winner cap.  These statements reduce the unresolved case
to a joint core in which every known cover, factor elimination, and
submodular escape is expensive; the remaining geometric obligation is a
lower bound on shortest direct mass plus shortest cone-volume mass.

We then prove that this geometric obligation would still not complete the
proposed empirical line scan.  Hhan's weighted shell estimate suppresses the
non-colliding theta tail on a homothetic facet core only when the retained
area is exponentially smaller than required.  More decisively, at every
tail-suppressed point whose theta response lies above the sampling noise
floor, the tangential displacement from the midpoint is confined to an
explicit Euclidean ball.  Even if all \(2^{n/2+o(n)}\) samples are spent on
one query, that ball contains at most \(2^{-(1.0289+o(1))n}\) of a shortest
facet.  Retaining \(2^{-(0.0264828+o(1))n}\) would require sample exponent at
least \(0.61169\ldots\).  Thus polynomial weighted-tail suppression and a
sufficient empirical theta-response floor cannot coexist on enough of the
robust facet core within the half-exponential budget.

Independently of that hypothesis, we transfer the moving-stabilizer
spherical-code bound of \cite{OpenAI26} through Hhan's shell estimates.  Its
one-row subfamily gives exponent $0.602470$, subject to the cited theorem and
a one-dimensional numerical enclosure.  Thus the note contains both a
checkable update of Hhan's parameter ledger and a conditional route to the
half exponent.

Conditioned on a replacement shortest-parity basin hypothesis for an
adapted law with directly observable endpoints, we obtain a randomized
classical algorithm for search-$\SVP$ in expected time and worst-case space
$2^{0.5n+o(n)}$, up to polynomial factors in the input length.  More
generally, an attraction probability $2^{-(\delta+o(1))n}$ suffices for every
$\delta<0.02648$.  Thus the note does not prove an unconditional
half-exponent algorithm: it proves the analytic reduction, refutes the first
universal candidate, and proves that the natural log-theta line detector
cannot repair that candidate at the desired exponent.  The surviving target
must control response-weighted shortest collisions or use a different
nonlocal estimator; direct-plus-cone mass alone is insufficient.
\end{abstract}

\medskip
\noindent\textbf{AI use disclosure.}
The technical route, finite experiments, and initial manuscript draft were
developed with substantial assistance from OpenAI Codex.  Every theorem
whose proof is complete is separated from the explicitly labeled geometric
hypothesis.  The human author of any circulated version is responsible for
checking the arguments, replacing the author placeholder, and deciding which
claims are suitable for dissemination.

\tableofcontents

\section{Introduction}

Let $L=L(B)\subset\R^n$ be a full-rank lattice.  The search shortest vector
problem asks for a nonzero $v\in L$ of norm $\lambda_1(L)$.  The best
previously proved general classical algorithm ran in $2^{n+o(n)}$ time and
space through discrete Gaussian sampling \cite{ADRS15}.  Hhan's recent
mid-point Hessian algorithm \cite{Hhan26} reduces the classical running time
to $2^{0.6039n+o(n)}$ while using $2^{n/2+o(n)}$ space.

The central observation of \cite{Hhan26} is local.  If $v$ is shortest and
$s=\Theta(\lambda_1(L)/\sqrt n)$ is chosen correctly, the Hessian of the
periodic Gaussian at $v/2$ is, up to a scalar identity matrix and small
operator error, a positive multiple of $vv^T$.  Its leading eigenvector is
therefore inverse-polynomially close to $v/\norm v$.  The preprocessing
$\BDD$ oracle of Aggarwal, Chen, Kumar, and Shen \cite{ACKS25} then recovers
$v$ in subexponential query time.

The expensive part is finding the parity class of $v$ in $L/2L$.  Hhan uses
matrix-valued Walsh--Hadamard transforms, random affine sublattice cosets, and
importance sampling.  The present note asks what can be obtained from the
same samples without first committing to the discrete set $L/2L$.
Hhan's Remark 1.2 already anticipates that the new code and packing bounds of
\cite{OpenAI26} lower his constant.  Section~\ref{sec:spherical-transfer}
gives a conservative, reproducible transfer using only the explicit one-row
spherical subfamily; the continuous route begins after that direct update.

\subsection{Technical overview}

For $s>0$, put $\rho_s(x)=\exp(-\pi\norm{x}^2/s^2)$ and
\[
 F_s(z)=\frac{\rho_s(L+z)}{\rho_s(L)}
      =\E_{X\sim D_{L^*,1/s}}\!\left[e^{2\pi i\ip Xz}\right].
\]
Differentiation gives the dual representation
\[
 \nabla^2 F_s(z)
   =-4\pi^2\E\left[XX^T\cos(2\pi\ip Xz)\right].
\]
The radial identity component is irrelevant to the eigenvectors.  We remove
it at the sample level and work with
\[
 \cT_s(z):=\nabla^2F_s(z)-\frac{\tr(\nabla^2F_s(z))}{n}I_n
 =-4\pi^2\E\left[Q(X)\cos(2\pi\ip Xz)\right],
 \quad
 Q(X):=XX^T-\frac{\norm X^2}{n}I_n.                 \tag{1}
\]
Every target sample can be reused at an arbitrary $z\in\Torus$, and
differentiating (1) supplies continuous gradients and Hessians of the matrix
field.

There are two exact structural facts.  First, every half-grid point
$B\theta/2$, $\theta\in\F_2^n$, is stationary for every entry of $\cT_s$ at
every scale.  Second, if $\theta$ is the coefficient parity of a shortest
vector $v$, then $B\theta/2=v/2\bmod L$.  Hence a continuous method need not
transport a critical point between scales.  It only needs to enter the
rounding cell of the correct parity; it can then snap to the exact midpoint
and invoke Hhan's target-scale recovery theorem.

To obtain a smooth vector field even at eigenvalue collisions, define
\[
 \Psi_{t,d,\gamma}(z)
 :=\frac{\sigma_{t,d}}{\gamma}
   \log\tr\exp\!\left(\frac{\gamma\cT_t(z)}{\sigma_{t,d}}\right),       \tag{2}
\]
where $\gamma=\poly(n)$ and $\sigma_{t,d}$ is an explicit midpoint signal
scale.  The gradient of (2) is analytic everywhere.  Haar initialization on
the torus followed by a fixed discretization of this gradient flow and exact
half-grid snapping is our seed law.

The scale used for ascent is below the direct sampling threshold.  We draw
samples at a target parameter $r>t_0$ and narrow them to $t<r$ using weights
bounded by one.  For
\[
 r=0.23675858,\qquad t=0.16,
\]
the auxiliary second-moment parameter and importance exponent are
\[
 q=\frac{tr}{2r-t}=0.1208271\ldots>\frac{t_0}{2},\qquad
 \iota(t,r)=\frac12\log_2\frac{r^2}{t(2r-t)}=0.0801086\ldots.          \tag{3}
\]
Thus the uniformly accurate smooth oracle uses
\[
 2^{(2t+\iota(t,r)+o(1))n}=2^{(0.4001086+o(1))n}                     \tag{4}
\]
samples.  The target midpoint estimator uses
$2^{(2r+o(1))n}=2^{(0.4735172+o(1))n}$ samples.  Both fit below the
$2^{n/2+o(n)}$ discrete-Gaussian and preprocessing-$\BDD$ cost.

\subsection{Contributions and logical status}

The following components are proved in this note.
\begin{enumerate}[leftmargin=2em]
\item The moving-stabilizer spherical-code bound substitutes monotonically
into Hhan's shell argument; its one-row specialization gives exponent below
$0.602470$ after the displayed numerical enclosure.
\item The full-torus traceless Hessian and all fixed-order physical
derivatives are expectations of explicit matrix trigonometric polynomials.
\item Every half-grid point is stationary at every scale, and rounding an
endpoint to the correct parity produces exactly the shortest midpoint modulo
$L$.
\item Narrowing from $r$ to $t$ uses weights in $[0,1]$.  Its exact relative
second moment has exponent $\iota(t,r)$ whenever the auxiliary parameter
$q=tr/(2r-t)$ lies above $t_0/2$.
\item One importance-weighted sample panel uniformly estimates $\cT_t$ and
its first two derivatives over the entire torus.  The uniform event permits
adaptive queries without a pathwise union bound.
\item Log-trace-exponential smoothing transfers the matrix estimates to a
gap-free scalar ascent field.
\item A robust shortest-parity basin lower bound implies a
$2^{n/2+o(n)}$-time and space algorithm by exact snapping, target-scale
midpoint estimation, preprocessing $\BDD$, and verification.
\item The universal Haar-seeded soft-Hessian law does not have the required
basin mass.  An anisotropic rectangular family gives success probability at
most $2^{-n+1}$ for every polynomial-step, polynomial-description schedule.
\item Short-dual syndrome equations, a coefficient Hamming bound, and a
stable spectral-tube lifting decomposition give unconditional
lattice-adapted parity covers.  Whenever their displayed entropy is $o(n)$,
they instantiate the half-exponential reduction without an ascent
hypothesis.
\item Bounded coefficient energy has an exact weighted-junction-width
factorization.  A displayed elimination order of width below the query
allowance solves that input by min--sum elimination and backward preimage.
\item If a basis extends to an obtuse superbase, a discrete coarea identity
reduces exact $\SVP$ to a global minimum cut.  In particular the previously
used simplex-cancellation family is polynomial-time solvable.
\item Infinitesimal causal transport through the Voronoi tiling has an exact
facet-area law.  Its cone-volume weights obey an isotropic matrix
conservation identity, and shortest facets have explicit robust cores.
\item The centered Hessian of the log theta function is exactly a lattice-site
posterior covariance.  At a two-site facet collision it exposes $vv^T/4$
independently of tangential position, reducing line detection to explicit
tail and relative-response estimates.
\item Those two estimates obey an incompatible exponent ledger.  Weighted
shell suppression on a homothetic core forces exponential area loss, while
the empirical response floor confines accessible points to a ball occupying
at most $2^{-(1.0289+o(1))n}$ of a shortest facet at the full
$2^{n/2+o(n)}$ sample budget.  Retaining exponent $0.0264828$ would require
sample exponent at least $0.61169\ldots$.
\item A noncollinear length gap gives an explicit winner cap.  The query
ledger closes from this clause alone once the gap ratio exceeds
$5.2670066\ldots$.
\end{enumerate}

The counterexample separates the proved reduction from its failed first
instantiation.  Hypothesis~\ref{hyp:basin} records the condition needed by
the reduction for the stated law; Theorem~\ref{thm:obstruction} proves that
condition false.  The lattice-adapted covers and exact factor transports
below remove this obstruction on a broad solvable class.  A universal
half-exponential algorithm therefore requires more than a direct-plus-cone
mass theorem on the residual joint core: it needs response-weighted collision
mass or a different observable whose empirical error is not divided by the
local theta value.

\section{Preliminaries}

We follow the notation and normalization of \cite{Hhan26}.  All logarithms
are natural unless a base is shown.  The dimension tends to infinity, and
$o(1)$ terms are uniform over fixed parameters in the stated ranges.  The
input basis is rational and has polynomial bit length.

\subsection{Lattices, Gaussian parameters, and decoding}

For a nonsingular $B\in\R^{n\times n}$, let $L=B\Z^n$ and
$L^*=B^{-T}\Z^n$.  Write $\lambda=\lambda_1(L)$.  As in
\cite{Hhan26}, LLL reduction and a polynomial scale grid permit us to assume
that a value $d$ is available with
\[
 \lambda\le d\le(1+1/n)\lambda.                                    \tag{5}
\]
Let $\beta_0=2^{0.4014\ldots}$ be the lattice-point counting constant used in
\cite{Hhan26}, and define
\[
 \xi_a(d):=\sqrt{\frac{4na\ln2}{\pi d^2}},
 \qquad
 t_0:=\frac{\beta_0^2}{4e\ln2}=0.23147\ldots.                       \tag{6}
\]
We use two black-box results exactly as in Hhan's algorithm.

\begin{theorem}[Above-smoothing DGS, \cite{ADRS15,Hhan26}]
If $a>t_0$ is fixed, one call at parameter $\xi_a(d)$ outputs up to
$2^{n/2}$ samples jointly negligibly close to independent samples from
$D_{L^*,\xi_a(d)}$ in time and space $2^{n/2+o(n)}$.
\end{theorem}

\begin{theorem}[Preprocessing BDD, \cite{ACKS25,Hhan26}]
There is randomized preprocessing using time and space $2^{n/2+o(n)}$ such
that, after successful preprocessing, inverse-polynomial $\BDD$ queries take
$2^{o(n)}$ time.  Outputs on targets outside the decoding promise may be
arbitrary, but exact lattice membership and norm can be verified.
\end{theorem}

The last verification clause is useful: every wrong parity or unstable
ascent path is harmless except for its running time.

\subsection{The target midpoint recovery lemma}

For $a>0$, abbreviate
\[
 F_a(z):=F_{1/\xi_a(d)}(z),\qquad H_a(z):=\nabla^2F_a(z).
\]
For $\theta\in\F_2^n$, choose its lift in $\{0,1\}^n$.  Hhan defines
$G_{a,d}(\theta)=H_a(B\theta/2)$ and adds a scalar identity to obtain a
positive semidefinite matrix $A_{a,d}(\theta)$.  The following packages
Lemmas 3.4--3.6 of \cite{Hhan26} in the form used here.

\begin{theorem}[Midpoint recovery, \cite{Hhan26}]                  \label{thm:hhan-endpoint}
Fix $t_0<r<1/4$ and suppose (5) holds.  Let $v$ be shortest and
$v\in B\theta+2L$.  There is an explicit
\[
 \mu_{r,d}=\xi_r^4\lambda^2\,2^{-rn+O(1)}                        \tag{7}
\]
such that, after addition of a scalar identity,
\[
 H_r(B\theta/2)=\mu_{r,d}\frac{vv^T}{\lambda^2}+E+aI_n,
 \qquad \norm E_{\op}\le \mu_{r,d}2^{-\Omega(n)}.                 \tag{8}
\]
If an estimator of $H_r(B\theta/2)$ has operator error
$O(\mu_{r,d}/n)$, then a leading eigenvector is
$O(n^{-1/2})$-close, up to sign, to $v/\lambda$.  A preprocessing $\BDD$
query at the corresponding scaled direction returns $\pm v$ for all
sufficiently large $n$.
\end{theorem}

Because adding or subtracting a scalar identity does not change eigenspaces,
Theorem~\ref{thm:hhan-endpoint} remains valid for the traceless Hessian used
below.

\section{A direct spherical-bound improvement}                   \label{sec:spherical-transfer}

Hhan's lattice-point estimates use the cap-optimized
Kabatianskii--Levenshtein spherical-code exponent in two places: the global
constant $\log_2\beta$ and the shell function $K_2$.  The binary-code theorem
of \cite{OpenAI26} does not itself apply to these Euclidean directions, but
the spherical theorem in the same work does.  We record the substitution
because it yields a lower constant without the geometric hypothesis used
later.

For $u\ge0$, put
\[
 H_{\rm sph}(u)=(1+u)\log_2(1+u)-u\log_2u,
 \qquad H_{\rm sph}(0)=0,
\]
and for $a>b\ge0$ put
\[
 \Gamma_{\rm row}(a,b)
 =\frac{(a-b)(1+a+b)}{(1+2a)\sqrt{a(1+a)}}.                       \tag{S1}
\]
The moving one-row theorem of \cite[Chapter 2]{OpenAI26} gives the valid
whole-sphere exponent
\[
 B_{\rm row}^{\rm dir}(s)
 :=\inf_{\substack{a>b\ge0\\2\Gamma_{\rm row}(a,b)\ge s}}
     \bigl(H_{\rm sph}(a)-H_{\rm sph}(b)\bigr),                  \tag{S2}
\]
and its cap optimization can only decrease this value.  Define
\begin{align}
 b_{\rm row}&:=B_{\rm row}^{\rm dir}(1/2),
 &t_{0,{\rm row}}&:=\frac{2^{2b_{\rm row}}}{4e\ln2},             \tag{S3}\\
 K_{\rm row}(x)&:=
 \begin{cases}
 0,&1\le x\le\sqrt2,\\
 B_{\rm row}^{\rm dir}(1-2/x^2),&x>\sqrt2,
 \end{cases}                                                     \tag{S4}\\
 g_{2,{\rm row}}(R)&:=1-\sup_{x\ge1}
 \left(\min\{b_{\rm row}+\log_2x,1+K_{\rm row}(x)\}
       -\frac{Rx^2}{2}\right).                                  \tag{S5}
\end{align}
For $g_{\rm row}(r)=\tfrac12\log_2(r/t_{0,{\rm row}})$ and
$\iota(r,R)=\tfrac12\log_2(R^2/[r(2R-r)])$, Hhan's classical exponent
ledger becomes
\begin{align}
 \chi(r,R)&=\min\{g_{2,{\rm row}}(R),\tfrac12+g_{\rm row}(r),
 1+\min(g_{\rm row}(r),2g_{\rm row}(r))-2r,\tfrac12\},          \tag{S6}\\
 c(r,R)&=\max\{\tfrac12,2r+\iota(r,R),1-\chi(r,R)\}.            \tag{S7}
\end{align}

\begin{theorem}[Moving-row transfer, numerical form]              \label{thm:moving-row}
Assume the spherical-code theorem of \cite[Chapter 2]{OpenAI26} and the
numerical enclosure (S8).  Hhan's
classical mid-point Hessian algorithm has time complexity
\[
 2^{0.602470n+o(n)}
\]
and space complexity $2^{n/2+o(n)}$.  More precisely, direct substitution of
\[
 r=0.22189254,\qquad R=0.39935521
\]
in (S2)--(S7) gives
\[
 b_{\rm row}=0.39730560\ldots,\quad
 t_{0,{\rm row}}=0.23015577\ldots,\quad
 \chi=0.39753095\ldots,
\]
\[
 2r+\iota(r,R)=0.60246910\ldots,
 \qquad 1-\chi=0.60246905\ldots.                                \tag{S8}
\]
\end{theorem}

\begin{proof}
Replace the spherical-code exponent in the shell count underlying Hhan's
Lemma 2.1 by (S2).  This changes $\log_2\beta$ to $b_{\rm row}$ and therefore
changes the Gaussian threshold to (S3).  In Hhan's refined shell bound,
replace the occurrence of the classical exponent at angle $1-2/x^2$ by
(S4).  Both replacements are pointwise upper bounds, so every subsequent
Gaussian-mass, affine-coset, importance-sampling, and concentration inequality
is unchanged except for the scalar functions in (S3)--(S7).  The original
optimization proof therefore gives exponent $c(r,R)$.

The two variational problems in (S2) and (S5) are one-dimensional after
solving $2\Gamma_{\rm row}(a,b)=s$ for $a$.  Evaluating them at the displayed
parameters gives (S8); the gap to $0.602470$ is more than $9\cdot10^{-7}$.
The nested formulas are evaluated directly rather than replaced by the
binary-code exponent.  Hence $c(r,R)<0.602470$.  Hhan's space argument is
unaffected and remains $2^{n/2+o(n)}$.
\end{proof}

\begin{remark}
Theorem~\ref{thm:moving-row} is logically independent of the continuous
basin hypothesis below.  It is a direct improvement of Hhan's proved
algorithm if the cited spherical-code theorem and the numerical enclosure in
(S8) are accepted.  The full stabilizer hierarchy of \cite{OpenAI26} is
pointwise no worse than the one-row family, but no additional numerical
improvement from that hierarchy is claimed here.
\end{remark}

\section{The continuous periodic-Hessian oracle}

\subsection{Traceless Fourier representation}

For $X\in\R^n$, define the traceless quadratic matrix
\[
 Q(X):=XX^T-\frac{\norm X^2}{n}I_n.                                \tag{9}
\]
For a scale $a>0$, define
\[
 \cT_a(z):=H_a(z)-\frac{\tr H_a(z)}{n}I_n.                         \tag{10}
\]

\begin{lemma}[Full-torus oracle]                                  \label{lem:full-torus}
For every $z\in\R^n$,
\[
 \cT_a(z)=-4\pi^2\E_{X\sim D_{L^*,\xi_a}}
       \left[Q(X)\cos(2\pi\ip Xz)\right].                        \tag{11}
\]
For every fixed integer $j\ge0$ and directions
$h_1,\ldots,h_j\in\R^n$,
\begin{align}
 D^j\cT_a(z)[h_1,\ldots,h_j]
 ={}&-4\pi^2(2\pi)^j\E\left[
 Q(X)\prod_{m=1}^j\ip X{h_m}
 \cos\!\left(2\pi\ip Xz+\frac{j\pi}{2}\right)
 \right].                                                        \tag{12}
\end{align}
In particular, one sample panel gives unbiased estimators at arbitrary,
adaptively selected torus points.
\end{lemma}

\begin{proof}
Poisson summation gives
$F_a(z)=\E[e^{2\pi i\ip Xz}]$.  The distribution is symmetric under
$X\mapsto-X$, so the imaginary part vanishes.  Differentiating twice and
subtracting the normalized trace gives (11).  Further directional
differentiation uses
$d^j\cos u/du^j=\cos(u+j\pi/2)$ and proves (12).  The final statement follows
because $z$ appears only in the trigonometric factor; the sample distribution
does not depend on the query.
\end{proof}

The traceless projection is not cosmetic.  It removes the radial identity
fluctuation from every individual sample while preserving every eigenvector
and eigengap of $H_a(z)$.

\subsection{Half-grid stationarity and exact parity snapping}

The next fact does not require a Gaussian estimate or a spectral gap.

\begin{lemma}[Scale-uniform half-grid stationarity]                \label{lem:stationary}
For every $a>0$ and $\theta\in\Z^n$,
\[
 D\cT_a(B\theta/2)=0.                                             \tag{13}
\]
Consequently, every differentiable spectral function of $\cT_a$, including
the soft maximum in Section~\ref{sec:soft}, is stationary at
$B\theta/2$.
\end{lemma}

\begin{proof}
Write $X=B^{-T}k$ with $k\in\Z^n$.  At $z=B\theta/2$,
$2\pi\ip Xz=\pi k^T\theta$, and every sine factor in the first derivative of
(11) vanishes.  The chain rule proves the second claim.
\end{proof}

\begin{lemma}[Parity snapping]                                    \label{lem:snap}
Let $c\in\R^n$ and suppose there is a $\theta\in\Z^n$ with
\[
 \norm{2c-\theta}_\infty<\frac12.                                 \tag{14}
\]
Then coordinate rounding of $2c$ recovers $\theta\bmod2$.  If a shortest
vector $v$ satisfies $v=B\theta+2w$ for some $w\in L$, then
\[
 B\theta/2=v/2\pmod L.                                            \tag{15}
\]
Therefore $\cT_r(B\theta/2)=\cT_r(v/2)$.
\end{lemma}

\begin{proof}
Condition (14) gives unique coordinate rounding.  From
$v=B\theta+2w$, divide by two to obtain $B\theta/2=v/2-w$.  Periodicity of
$F_r$ and its derivatives gives the last claim.
\end{proof}

Thus a smooth-scale critical point is not transported to the target scale.
Its parity is transported algebraically by one exact rounding operation.

\section{Narrowing importance sampling}

Fix a source parameter $r>t_0$ and a smaller ascent parameter $t<r$.  Put
$\xi_a=\xi_a(d)$ and let $P_a$ denote $D_{L^*,\xi_a}$.  Since
$\xi_t<\xi_r$, define
\[
 u_t(x):=\frac{\rho_{\xi_t}(x)}{\rho_{\xi_r}(x)}
 =\exp\!\left[-\pi\norm x^2
       \left(\frac1{\xi_t^2}-\frac1{\xi_r^2}\right)\right].       \tag{16}
\]
The elementary but useful point is
\[
 0<u_t(x)\le1.                                                     \tag{17}
\]
There is no upper-tail problem in the importance weight; the loss comes from
its small normalizing mean.

\begin{lemma}[Exact narrowing identity]                            \label{lem:importance}
For every integrable matrix-valued $A$,
\[
 \E_t[A(X)]
 =\frac{\E_r[u_t(X)A(X)]}{\E_r[u_t(X)]}.                           \tag{18}
\]
Let
\[
 q=q(t,r):=\frac{tr}{2r-t},\qquad
 \iota(t,r):=\frac12\log_2\frac{r^2}{t(2r-t)}.                    \tag{19}
\]
Then
\[
 \frac{\E_r[u_t(X)^2]}{\E_r[u_t(X)]^2}
 =2^{\iota(t,r)n}
 \frac{\rho_{1/\xi_q}(L)\rho_{1/\xi_r}(L)}
      {\rho_{1/\xi_t}(L)^2}.                                     \tag{20}
\]
If $q>t_0/2$, the final theta-mass ratio is $1+2^{-\Omega(n)}$ and
the relative second moment is $2^{(\iota(t,r)+o(1))n}$.
\end{lemma}

\begin{proof}
Expanding both expectations in (18) cancels $\rho_{\xi_r}$ and gives the
$P_t$ expectation after division by the mass ratio.  Moreover,
\[
 \frac{\rho_{\xi_t}(x)^2}{\rho_{\xi_r}(x)}=\rho_{\xi_q}(x)
\]
because $1/q=2/t-1/r$.  It follows that the left side of (20) is
\[
 \frac{\rho_{\xi_q}(L^*)\rho_{\xi_r}(L^*)}
      {\rho_{\xi_t}(L^*)^2}.
\]
Poisson summation changes the power-of-scale factor into
$(qr/t^2)^{n/2}=2^{\iota(t,r)n}$ and leaves the displayed primal theta-mass
ratio.  Hhan's Gaussian mass bound implies
$\rho_{1/\xi_a}(L)=1+2^{-\Omega(n)}$ for every fixed $a>t_0/2$.
Since $q<t<r$, the assumption on $q$ proves the last claim.
\end{proof}

\section{A uniform empirical derivative oracle}

We now show that adaptivity does not require fresh samples.  All derivatives
below are physical derivatives on $\R^n$, not coefficient derivatives; this
avoids introducing the condition number of the input basis into the random
summands.

Let
\[
 \overline\mu_{t,d}:=\xi_t^4d^2\,2^{-tn}.                         \tag{21}
\]
For the correct length guess, this is the coefficient scale of the
shortest-pair contribution, up to polynomial and constant factors.  It does
not assert that this rank-one contribution dominates the full Hessian when
$t<t_0$; that geometric issue is precisely what the basin hypothesis below
must control.

Given independent $X_1,\ldots,X_M\sim P_r$, define the ratio estimator
\[
 \what\cT_t(z):=
 -4\pi^2\frac{\sum_{i=1}^M u_t(X_i)Q(X_i)
                    \cos(2\pi\ip{X_i}{z})\mathbf1_{\cG_i}}
                   {\sum_{i=1}^M u_t(X_i)\mathbf1_{\cG_i}},        \tag{22}
\]
where $\cG_i$ is the event
$\norm{X_i}\le C\xi_r n$ for a sufficiently large fixed $C$.
Its derivative estimators are obtained by differentiating the
trigonometric factor in (22).

\begin{theorem}[Uniform narrowing oracle]                          \label{thm:uniform}
Fix constants $r,t$ with
\[
 t_0<r<1/4,\qquad 0<t<r,\qquad q(t,r)>t_0/2.                       \tag{23}
\]
For every fixed $A>0$, there is a fixed $C_A$ such that, with
\[
 M=(n+\mathrm{size}(B))^{C_A}\,2^{(2t+\iota(t,r))n},              \tag{24}
\]
the following holds with probability at least
$1-\exp(-n^2)$: simultaneously for every $z\in\Torus$ and
$j\in\{0,1,2\}$,
\[
 \sup_{\norm{h_1},\ldots,\norm{h_j}\le1}
 \norm{D^j\what\cT_t(z)[h_1,\ldots,h_j]
      -D^j\cT_t(z)[h_1,\ldots,h_j]}_{\op}
 \le \frac{\overline\mu_{t,d}\xi_t^j}{n^A}.                     \tag{25}
\]
The same sample panel supports an arbitrary adaptive sequence of queries.
For ideal independent samples the failure probability is at most
$\exp(-n^2)$.  For the actual DGS output it is at most $2^{-\Omega(n)}$ after
the statistical-distance loss.
\end{theorem}

\begin{proof}
We expose the exponent calculation and then account for uniformity.
For a fixed $z$, directions $h_m$, and matrix entry, the numerator summand
for the $j$th derivative is bounded on $\cG_i$ by
\[
 (C'\xi_r n)^{2+j}.                                               \tag{26}
\]
The omitted Gaussian tail contributes at most
$\overline\mu_{t,d}\xi_t^j/n^{A+2}$ after increasing $C$; this is the
standard discrete-Gaussian tail integration used in \cite{Hhan26}.

Let $m_t=\E_r[u_t(X)]$.  By Lemma~\ref{lem:importance},
\[
 \frac{\E_r[u_t^2]}{m_t^2}=2^{(\iota(t,r)+o(1))n}.                 \tag{27}
\]
Scalar Bernstein concentration applied to the bounded denominator gives
$\sum_i u_t(X_i)=(1+o(1))Mm_t$.  Matrix Bernstein applied to the centered
weighted numerator, together with (27), has the same variance as
$M2^{-(\iota+o(1))n}$ target samples.  The target signal scale is
$2^{-tn}$ up to polynomial factors, so inverse-polynomial relative accuracy
requires $2^{(2t+\iota)n}$ samples.  The polynomial $n^{C_A}$ makes the
failure probability at a fixed query smaller than
$\exp(-n^{C'})$ for an arbitrarily large fixed $C'$.

It remains to make the event uniform.  On $\cG_i$, (12) bounds the next
physical derivative of every summand by a number whose logarithm is
polynomial in $n$ and the input bit length.  A rational fundamental
parallelepiped of $L$ has diameter at most exponential in that bit length.
Consequently it has a net fine enough for (25) whose logarithmic cardinality
is polynomial in $n$ and the input bit length.  Apply the fixed-query bound
to the net, to a fixed operator-norm net on the unit sphere, and to
$j=0,1,2$; choose $C'$ above the resulting polynomial; and use the next
derivative bound to extend from the net to the whole torus.  This proves
(25).  Since the event already quantifies over every $z$, adaptively chosen
queries require no additional union bound.  The statistical distance in the
DGS theorem adds at most $2^{-\Omega(n)}$ to the ideal-sample failure
probability.
\end{proof}

\begin{remark}
The theorem is a statement about oracle accuracy, not automatic trajectory
shadowing.  Near a separatrix, an arbitrarily small perturbation can change
the terminal basin.  The robust basin hypothesis in Section~\ref{sec:basin}
quantifies the dynamical margin that is additionally required.
\end{remark}

\section{A gap-free soft Hessian objective}                        \label{sec:soft}

The map $T\mapsto\lambda_{\max}(T)$ is not differentiable when the leading
eigenvalue is multiple.  This is inconvenient precisely at symmetric
lattices.  For $\gamma>0$ and $\sigma=\overline\mu_{t,d}$, define
\[
 \psi_{\gamma,\sigma}(T)
 :=\frac{\sigma}{\gamma}\log\tr e^{\gamma T/\sigma},
 \qquad
 \Psi_{t,d,\gamma}(z):=\psi_{\gamma,\sigma}(\cT_t(z)).             \tag{28}
\]

\begin{lemma}[Soft spectral calculus]                              \label{lem:soft}
For every symmetric $T$,
\[
 \lambda_{\max}(T)
 \le\psi_{\gamma,\sigma}(T)
 \le\lambda_{\max}(T)+\frac{\sigma\log n}{\gamma}.              \tag{29}
\]
If
\[
 P_T:=\frac{e^{\gamma T/\sigma}}{\tr e^{\gamma T/\sigma}},
\]
then
\[
 D\psi_{\gamma,\sigma}(T)[E]=\tr(P_TE).                           \tag{30}
\]
Every derivative of $\psi$ is analytic for every symmetric $T$, and the
operator norm of its second Fr\'echet derivative is at most
$O(\gamma/\sigma)$.  Therefore Theorem~\ref{thm:uniform}, with a larger
polynomial factor, uniformly estimates $\Psi$ and $\nabla\Psi$ to
inverse-polynomial absolute error without any eigengap assumption.
\end{lemma}

\begin{proof}
If $\lambda_1\ge\cdots\ge\lambda_n$ are the eigenvalues, then
$e^{\gamma\lambda_1/\sigma}\le\sum_i e^{\gamma\lambda_i/\sigma}
\le ne^{\gamma\lambda_1/\sigma}$, proving (29).  The Fr\'echet derivative of
$\log\tr e^T$ gives (30).  Matrix exponentiation and trace are analytic, and
the trace exponential is strictly positive.  The derivative in (30) has
operator norm at most one, while the covariance formula for its derivative
is bounded by $O(\gamma/\sigma)$.  The physical first derivatives of
$\cT_t$ have polynomial norm by (12) and the Gaussian moment bounds.
Consequently an error $\sigma/n^A$ in $\cT_t$ and its first derivative
induces inverse-polynomial absolute error in $\Psi$ and its gradient after
increasing $A$.  This is exactly the uniform event of
Theorem~\ref{thm:uniform}.
\end{proof}

When a polynomial eigengap is present and $\gamma=\poly(n)$ is sufficiently
large, $P_T$ is exponentially close to the projector onto the leading
eigenvector.  When no gap is present, (30) remains a canonical convex mixture
of the tied eigenspaces rather than an undefined branch choice.

\section{The samplable seed law}

Fix $t=0.16$, $\gamma=n^2$, a polynomial step count $S=n^{C_0}$, and a
deterministic step schedule $\eta_0,\ldots,\eta_{S-1}$ of polynomial bit
complexity.  Let $\what\Psi$ be the empirical soft objective from the sample
panel in Theorem~\ref{thm:uniform}.  Define a bounded vector field
\[
 \what V(z):=
 \frac{d^4\nabla\what\Psi(z)}
      {1+d^3\norm{\nabla\what\Psi(z)}}.                            \tag{31}
\]
The normalization is globally Lipschitz and keeps the step bounded; any fixed
smooth normalization with polynomial Lipschitz bounds could be used instead.
The powers of $d$ make the field scale-equivariant: replacing $L$ by
$\alpha L$ replaces $V$ by $\alpha V$.

\medskip
\noindent\textbf{Soft-Hessian seed law $\widehat\nu_{L,d}$.}
\begin{enumerate}[leftmargin=2.2em,label=\arabic*.]
\item Draw $U\sim\mathrm{Unif}([0,1)^n)$ and put $Z_0=BU\bmod L$.
\item For $s=0,\ldots,S-1$, set
\[
 Z_{s+1}=Z_s+\eta_s\what V(Z_s)\pmod L.
\]
\item Compute $c=B^{-1}Z_S$ modulo $\Z^n$, set
$\theta=\operatorname{round}(2c)\bmod2$, and output $B\theta/2\bmod L$.
\end{enumerate}

This is a samplable probability law using one target DGS panel.  It does not
presuppose the shortest vector or its parity.  Its exact counterpart
$\nu_{L,d}$ replaces $\what\Psi$ by $\Psi$.  All transcendental evaluations
are performed to inverse-polynomial precision; the robustness margin below
absorbs that deterministic approximation.

\section{The geometric existence clause}                          \label{sec:basin}

Let $\Theta_{\min}\subseteq\F_2^n$ be the set of parities represented by
shortest vectors.  We need robustness because uniform oracle accuracy cannot
stabilize a point lying exactly on a basin boundary.

\begin{definition}[Robust shortest-parity basin]
For $\varepsilon>0$, let $\cS_{L,d}(\varepsilon)$ be the set of Haar seeds
$Z_0$ with the following property: every vector field $\widetilde V$ whose
value is uniformly within $\varepsilon d$ of the exact normalized soft field
$V$ produces, under the fixed $S$-step schedule, a
terminal point satisfying (14) with a parity in $\Theta_{\min}$ and with
rounding margin at least $\varepsilon$.
Set
\[
 p_{L,d}(\varepsilon)
 :=\Haar_{\Torus}(\cS_{L,d}(\varepsilon)).                          \tag{32}
\]
\end{definition}

\begin{hypothesis}[Robust shortest-parity basin]                   \label{hyp:basin}
There is a universal computable function $\kappa(n)=o(n)$ such that, for every
full-rank lattice $L$, at one scale guess $d$ satisfying (5), there is an
inverse-polynomial $\varepsilon=n^{-C_1}$ for which
\[
 p_{L,d}(\varepsilon)\ge2^{-\kappa(n)}.                           \tag{33}
\]
The step count, smoothing parameter, and constants $C_0,C_1$ are universal.
\end{hypothesis}

Hypothesis~\ref{hyp:basin} is exactly the additional statement used by the
$2^{n/2+o(n)}$ reduction for this candidate law.  A stronger
uniform-attraction theorem would not be needed, nor would a bound on the
total number of local maxima.  However, Theorem~\ref{thm:obstruction} below
shows that the hypothesis itself is false for the universal Haar-seeded,
single-scale law.

\begin{proposition}[Empirical preservation of robust mass]        \label{prop:preserve}
Assume Hypothesis~\ref{hyp:basin}.  Increase the polynomial factor in
Theorem~\ref{thm:uniform} so that the normalized empirical vector field is
within $\varepsilon d$.  Then a sample from
$\what\nu_{L,d}$ snaps to $\Theta_{\min}$ with probability
$2^{-\kappa(n)}-2^{-\Omega(n)}$.
\end{proposition}

\begin{proof}
On the uniform event of Theorem~\ref{thm:uniform}, Lemma~\ref{lem:soft} and
the normalization in (31) put $\what V$ in the perturbation class quantified
in the preceding definition.  Every seed in
$\cS_{L,d}(\varepsilon)$ therefore snaps to
a shortest parity.  The exceptional oracle event for the actual DGS panel has
probability at most $2^{-\Omega(n)}$.
\end{proof}

\section{The conditional \texorpdfstring{$2^{n/2+o(n)}$}{half-exponential}
algorithm}

We can now state the algorithm in the same verified-output model used by
Hhan.

\medskip
\noindent\textbf{Algorithm 1: Continuous mid-point Hessian $\SVP$.}
\begin{enumerate}[leftmargin=2.2em,label=\arabic*.]
\item Apply LLL, construct the polynomial grid of guesses $d$, and construct
the preprocessing $n^{-1/3}$-$\BDD$ data.
\item Fix $r=0.23675858$ and $t=0.16$.  At each $d$, call above-smoothing
$\DGS$ on $L^*$ at $\xi_r(d)$ and retain enough samples for the smooth oracle
and target midpoint estimators.  Abort a scale returning too few samples.
\item Using the same panel and the weights (16), draw
$K=\Theta(2^{\kappa(n)})=2^{o(n)}$ independent outputs of the soft-Hessian
seed law $\what\nu_{L,d}$.
\item For every snapped parity $\theta$, estimate the target traceless
Hessian $\cT_r(B\theta/2)$ from $n^{C}2^{2rn}$ unweighted source samples.
Compute a leading eigenvector $q$ and query $\BDD$ at $dq$ and $-dq$.
\item Verify every output, and return the shortest verified nonzero lattice
vector over all guesses and starts.
\end{enumerate}

\begin{theorem}[Conditional main theorem]                          \label{thm:main}
Assume Hypothesis~\ref{hyp:basin}.  Algorithm 1 solves search-$\SVP$ with
probability at least $2/3$ in expected time and worst-case space
\[
 2^{n/2+o(n)},
\]
up to factors polynomial in the input length.
\end{theorem}

\begin{proof}
Choose the correct length guess $d$.  By
Proposition~\ref{prop:preserve}, one start snaps to a shortest parity with
probability at least $2^{-\kappa(n)}-2^{-\Omega(n)}$.  Taking
$K=\Theta(2^{\kappa(n)})=2^{o(n)}$ independent Haar starts gives constant
probability of at least one correct snap.  Condition on this event and call
its parity $\theta$.  By
Lemma~\ref{lem:snap}, the target point is exactly $v/2\bmod L$ for a shortest
vector $v$.  The target estimator with $n^C2^{2rn}$ samples has operator
error $O(\mu_{r,d}/n)$ by the same concentration argument as Hhan's Lemma
3.5.  Theorem~\ref{thm:hhan-endpoint} makes one of the $\BDD$ queries return
$\pm v$.  Exact verification prevents any false candidate from affecting
the output.  Standard independent repetitions raise the success probability
to $2/3$.

The $\BDD$ preprocessing and one above-smoothing DGS call cost
$2^{n/2+o(n)}$ time and space.  A smooth field/gradient query costs
\[
 2^{(2t+\iota(t,r)+o(1))n}=2^{(0.4001086+o(1))n}.
\]
There are $S K=2^{o(n)}$ such queries.  Each target midpoint query costs
\[
 2^{(2r+o(1))n}=2^{(0.4735172+o(1))n},
\]
and there are $K=2^{o(n)}$ of them.  All remaining linear algebra, wrapping,
rounding, verification, and $\BDD$ query work has a subexponential
multiplicative overhead.  Thus preprocessing dominates at exponent $1/2$.
The sample panel and $\BDD$ preprocessing data fit in
$2^{n/2+o(n)}$ space.
\end{proof}

The exponent ledger also gives a quantitatively weaker sufficient basin
bound.

\begin{corollary}[Exponential basin allowance]                     \label{cor:delta}
Suppose (33) is replaced by
\[
 p_{L,d}(n^{-C_1})\ge2^{-(\delta+o(1))n}.
\]
Algorithm 1 still runs in $2^{n/2+o(n)}$ time whenever
\[
 \delta<\min\left\{
 \frac12-(2t+\iota(t,r)),\frac12-2r
 \right\}=0.0264828\ldots.                                       \tag{34}
\]
\end{corollary}

\begin{proof}
Use $K=2^{(\delta+o(1))n}$ starts.  The smooth and target work exponents
become $2t+\iota+\delta$ and $2r+\delta$, respectively.  Both remain below
$1/2$ under (34).
\end{proof}

\section{Foundational consequences}

We collect several consequences that are implicit in the Fourier
representation but are not needed in Hhan's discrete Walsh formulation.

\subsection{Information lives on the whole torus}

The samples used for one parity Hessian are not tied to that parity.  They
encode the entire trigonometric matrix field (11).  Evaluating at a new point
does not require new randomness; it requires only another deterministic pass
over the same samples.  The same is true for every fixed derivative order.
This converts parity search from a sample-access problem into a geometric
query problem.

\subsection{Scale transport is needed only for the observable}

The narrowing identity transports the Hessian observable from $r$ to $t$.
It does not transport particles between affine cosets, nor does it move a
midpoint branch.  Lemma~\ref{lem:stationary} shows that the parity half-grid
is fixed at all temperatures, and Lemma~\ref{lem:snap} transfers the discrete
address exactly.  The final rank-one certificate is always evaluated at the
original target scale $r>t_0$.

\subsection{No eigengap is needed during search}

The target midpoint needs a leading eigendirection, and Hhan proves the
required gap there.  The search trajectory need not have one.  The soft
spectral density $P_T$ in (30) supplies a canonical direction through
collisions.  Thus the geometric hypothesis concerns basin mass and dynamical
margin, not a uniform eigengap over the torus.

\subsection{All false candidates are one-sided errors}

The algorithm never trusts the value of the soft objective or the output of
an out-of-promise $\BDD$ query as a certificate.  It verifies lattice
membership exactly and returns the shortest verified nonzero vector.  The
only role of the basin hypothesis is to guarantee that at least one query is
correct; it is not used to exclude incorrect queries.

\section{Finite evidence}

The experiments in this section are included only as checks of the formulas
and as evidence for the geometric target.  They are not used in any proof.

\subsection{Soft ascent on a cancellation family}

Consider a basis with Gram matrix having diagonal one and a common negative
off-diagonal chosen so that the all-ones coefficient vector has length
$0.97$.  In the tested dimensions it is the unique shortest pair, while each
basis column has length one.  At the rigorously controlled scale
$t=0.159919\ldots$, $\gamma=16$, Fourier cutoff two, and 256 scrambled Sobol
starts, the soft law gave the following results.

\begin{center}
\begin{tabular}{rrrrr}
\toprule
$n$ & shortest hits & hit fraction & distinct parities & median evaluations\\
\midrule
3 & 65 & 0.253906 & 7  & 14\\
4 & 40 & 0.156250 & 15 & 16\\
5 & 30 & 0.117188 & 31 & 18\\
6 & 25 & 0.097656 & 50 & 19\\
\bottomrule
\end{tabular}
\end{center}

The maximum half-grid error among terminal points was below $2.9\cdot10^{-6}$.
The number of discovered parity maxima is already large, but the shortest
parity retains visible basin mass on this small, well-conditioned family.
Theorem~\ref{thm:obstruction} shows why such evidence cannot be extrapolated
uniformly to anisotropic lattices.

\subsection{Finite-population oracle checks}

On explicitly enumerated dual cubes, the population identity (18) held to
floating-point precision, and every narrowing weight was at most one as
predicted by (17).  Independent-source experiments showed decreasing value
and gradient error with sample size.  No finite-dimensional measurement is
used to infer an asymptotic basin lower bound.

\section{Resolution of the universal-basin conjecture}

The universal Haar-seeded law has a product-lattice obstruction.  It is not
merely that the preceding estimates fail to prove a large basin: the basin
claimed in the former Conjecture~13.1 can be exponentially small.

\begin{theorem}[Anisotropic product obstruction]                  \label{thm:obstruction}
Fix $t=0.16$, $\gamma=n^2$, any polynomial step count, and any deterministic
step schedule of polynomial bit complexity as in the soft-Hessian seed law.
For all sufficiently large $n$, let
\[
 M=2^n,
 \qquad
 L_n=\operatorname{diag}(1,M,\ldots,M)\Z^n.                      \tag{35}
\]
For every correct guess $1\le d\le1+1/n$, the exact soft-Hessian seed law
reaches a shortest-vector parity with probability at most
\[
 2^{-n}+2^{-2n+2}+4n2^{-4n}\le2^{-n+1}.                          \tag{36}
\]
Consequently, for every $\varepsilon>0$,
\[
 p_{L_n,d}(\varepsilon)\le2^{-n+1}.                              \tag{37}
\]
Hypothesis~\ref{hyp:basin}, and hence the former universal-basin
Conjecture~13.1, is false.
\end{theorem}

\begin{proof}
Write $B=\operatorname{diag}(1,M,\ldots,M)$ and use coefficient coordinates
$z=Bu$ on the torus.  The only shortest vectors of $L_n$ are
$\pm e_1$, so the unique shortest parity is
$\theta_*=e_1\bmod2$.

Let $s=1/\xi_t(d)$.  Since $t$ is fixed and $d\in[1,1+1/n]$, there are
absolute constants $c_0,c_1>0$ such that
\[
 c_0n^{-1/2}\le s\le c_1n^{-1/2}.                               \tag{38}
\]
For $h>0$, define the one-dimensional periodic Gaussian
\[
 g_{h,s}(x):=\sum_{m\in\Z}e^{-\pi(x+mh)^2/s^2},
 \qquad f_{h,s}(x):=\frac{g_{h,s}(x)}{g_{h,s}(0)}.
\]
Separation of the lattice sum gives
\[
 F_t(z)=f_{1,s}(z_1)\prod_{j=2}^n f_{M,s}(z_j).                  \tag{39}
\]

We use the following elementary localization estimate.  For fixed
$a\in(0,1/4)$ and $0\le k\le3$, differentiation of each Gaussian summand
shows, uniformly for $0<s\le h$,
\[
 \left|f_{h,s}^{(k)}(x)\right|
 \le P_{k,a}(h+s^{-1})
 \exp\!\left(-c_a h^2/s^2\right)                                \tag{40}
\]
whenever $\operatorname{dist}(x,h\Z)\ge ah$.  Without the distance
assumption the same derivative is at most $P_k(h+s^{-1})$.  Here the $P$'s
are fixed polynomials and $c_a>0$.  Indeed, the $k$th derivative of
$e^{-\pi y^2/s^2}$ is a degree-$k$ polynomial in $y/s^2$ times the same
Gaussian.  Split the sum at a nearest translate of $h\Z$; the nearest term
has $|y|\ge ah$, and the two remaining tails are bounded by geometric
Gaussian tails.  Division is harmless because $g_{h,s}(0)\ge1$.  Finally,
$0<f_{h,s}\le1$: Poisson summation expresses its numerator as a Fourier
series with nonnegative coefficients, whose absolute value is at most its
value at zero.

Apply (40) with the constant $a/2$.  Suppose some long coordinate satisfies
\[
 \operatorname{dist}(u_j,\Z)\ge a/2,
 \qquad j\ge2.                                                   \tag{41}
\]
Then $\operatorname{dist}(z_j,M\Z)\ge aM/2$.  Every third partial derivative
of the product (39) contains either $f_{M,s}(z_j)$ or one of its first three
derivatives.  The other undifferentiated factors are at most one, and at most
three other factors are differentiated.  Equations (38)--(40) therefore
give a fixed polynomial $P$ and a constant $c>0$ such that
\[
 \max_{|\alpha|=3}|\partial^\alpha F_t(z)|
 \le P(n,M)e^{-cnM^2}.                                          \tag{42}
\]
The derivative of the traceless Hessian is a linear combination of these
third derivatives.  Moreover, by (30),
\[
 |\partial_j\Psi(z)|
 =|\tr(P_{\cT_t(z)}\,\partial_j\cT_t(z))|
 \le\norm{\partial_j\cT_t(z)}_{\op},                            \tag{43}
\]
because $P_{\cT_t(z)}$ is positive semidefinite with trace one.  Thus no
factor involving $1/\sigma_{t,d}$ is lost.  From (31), (42), and (43), after
enlarging $P$,
\[
 \norm{V(z)}\le P(n,M)e^{-cnM^2}                                \tag{44}
\]
whenever (41) holds.

Let $E_n=\sum_{s=0}^{S-1}|\eta_s|$.  On the input (35), polynomially many
step sizes of polynomial bit complexity satisfy
$E_n\le2^{q(n)}$ for some fixed polynomial $q$.  Since $M=2^n$, (44) implies
for all sufficiently large $n$ that
\[
 E_nP(n,M)e^{-cnM^2}<\tau_n,
 \qquad \tau_n:=2^{-4n}.                                       \tag{45}
\]

Set $a=1/8$ and define the exceptional active slabs
\[
 A:=\{u\in[0,1)^n:
       \operatorname{dist}(u_j,\Z)<a\text{ for every }j\ge2\}.
\]
Their Haar measure is
\[
 \Haar(A)=(2a)^{n-1}=2^{-2n+2}.                                 \tag{46}
\]
If $u_0\notin A$, choose a long coordinate satisfying (41).  The coefficient
displacement in one step is
$B^{-1}\eta_sV(Bu_s)$, whose norm is at most
$|\eta_s|\norm{V(Bu_s)}$ because $\norm{B^{-1}}_{\op}=1$.
Equations (44)--(45), applied inductively, show that the total coefficient
displacement is below $\tau_n<a/2$.  The witness coordinate therefore stays
at distance at least $a/2$ from $\Z$, so the localization estimate remains
valid at every step.

The coordinate rounding map $u\mapsto\operatorname{round}(2u)\bmod2$ changes
only across $1/4$ and $3/4$ in each circle coordinate.  Let $C$ be the set
of coefficient points within $\tau_n$ of one of these boundaries in at least
one coordinate.  A union bound gives
\[
 \Haar(C)\le4n\tau_n.                                           \tag{47}
\]
Outside $A\cup C$, the final parity equals the initial parity.  The initial
parity of a Haar point is uniform on $\F_2^n$.  Hence
\[
 \Prb[\text{final parity}=\theta_*]
 \le\Prb[\text{initial parity}=\theta_*]+\Haar(A)+\Haar(C),
\]
which is (36).

Finally, every seed in the robust set $\cS_{L_n,d}(\varepsilon)$ must
succeed for the exact vector field itself, since that field is one of the
allowed perturbations.  The robust set is therefore contained in the exact
success set.  This proves (37).  Since $n-1$ is not $o(n)$, no universal
$\kappa(n)=o(n)$ can satisfy (33).
\end{proof}

\subsection{An exact lattice-adapted parity sieve}

The product obstruction contains parity information that a Haar law ignores.
Let $v=Bz$ be shortest and put $\theta=z\bmod2$.  For
$k\in\Z^n$, write $y_k=B^{-T}k\in L^*$.

\begin{proposition}[Dual-syndrome and Hamming sieve]              \label{prop:sieve}
Suppose (5) holds.  If $d\norm{y_k}<1$, then
\[
 k^Tz=0,
 \qquad k^T\theta=0\pmod2.                                    \tag{48}
\]
Moreover, with
\[
 q(B,d):=\min\{n,\lfloor d^2\norm{B^{-1}}_{\op}^2\rfloor\},     \tag{49}
\]
every shortest parity is nonzero and has Hamming weight at most $q(B,d)$.
Consequently, for any matrix $K$ of verified rows satisfying (48), every
shortest parity lies in the explicitly defined set
\[
 \cP_{B,d,K}:=\{\theta\in\F_2^n\setminus\{0\}:
 K\theta=0,\ \operatorname{wt}(\theta)\le q(B,d)\},             \tag{50}
\]
whose cardinality is at most
\[
 Q_q:=\sum_{j=1}^{q(B,d)}\binom nj.                             \tag{51}
\]
It can be enumerated by filtering either the Hamming ball or the binary
kernel, in time polynomial in
\[
 E_{q,K}:=\min\{Q_q,2^{n-\operatorname{rank}_{\F_2}K}\}.         \tag{51a}
\]
\end{proposition}

\begin{proof}
The integer $k^Tz=\ip{y_k}{v}$ has absolute value at most
$d\norm{y_k}<1$, proving (48).  Also
\[
 \norm z\le\norm{B^{-1}}_{\op}\norm v
 \le d\norm{B^{-1}}_{\op}.
\]
Every odd coordinate of $z$ contributes at least one to $\norm z^2$, which
proves (49)--(51).  Finally, a shortest parity cannot be zero: if $z=2w$,
then the nonzero lattice vector $Bw=v/2$ is shorter.
\end{proof}

\begin{corollary}[Sparse-parity solvable class]                  \label{cor:sieve}
If $E_{q,K}=2^{o(n)}$, enumerating (50), applying the target midpoint
estimator and BDD endpoint to every retained parity, and verifying the
outputs solves $\SVP$ in $2^{n/2+o(n)}$ time and space.  More generally it
suffices that $E_{q,K}\le2^{(1/2+o(1))n}$ and
\[
 \frac1n\log_2|\cP_{B,d,K}|<0.0264828\ldots.                    \tag{52}
\]
In particular, $q=o(n)$ suffices.  If $q/n\to\rho$, it suffices that
$H_2(\rho)<0.0264828\ldots$, where
$H_2(\rho)=-\rho\log_2\rho-(1-\rho)\log_2(1-\rho)$, equivalently
$\rho<0.00264782\ldots$.
\end{corollary}

\begin{proof}
Every candidate midpoint query costs $2^{(2r+o(1))n}$ operations using the
common target sample panel, while the filtering work is polynomial in
$E_{q,K}$.  Combine this with (34); (51) has exponent
$H_2(\rho)+o(1)$ when $q/n\to\rho$.
\end{proof}

For the lattice (35), the rows $k=e_2,\ldots,e_n$ have
$\norm{y_k}=2^{-n}$.  They force every long parity bit to zero.  The only
remaining nonzero parity is $e_1$, so the counterexample to universal Haar
ascent is solved deterministically by the input geometry.

\subsection{A spectral support-chain cover}

The sieve can be vacuous when a shortest vector is a dense coefficient
cancellation.  The support-chain remedy is to transport the low-energy
coarse response and store only its sparse rounding detail.  Put $G=B^TB$.
Let $U\subset\R^n$ have dimension $k$ and orthogonal projector $P_U$, and
suppose every point of the coefficient ellipsoid
\[
 \cE_d:=\{z\in\R^n:z^TGz\le d^2\}                              \tag{53}
\]
satisfies
\[
 \norm z\le R,
 \qquad \norm{(I-P_U)z}\le C.                                 \tag{54}
\]
For the first $k$ Gram eigendirections one may take
$R=d/\sqrt{\lambda_1(G)}$ and
$C=d/\sqrt{\lambda_{k+1}(G)}$.

\begin{theorem}[Spectral-tube parity cover]                      \label{thm:tube}
Let $N=n(2\lceil R\rceil+2)$ and
$s=\min\{n,\lfloor4C^2\rfloor\}$, $m=\lceil C+1/2\rceil$.
Every shortest parity is contained in an explicitly enumerable cover of
cardinality at most
\[
 2^k\left(\sum_{j=0}^k\binom Nj\right)
 \left(\sum_{i=0}^s\binom ni(2m)^i\right).                     \tag{55}
\]
In particular the cover has size $2^{o(n)}$ whenever, uniformly in the
input model,
\[
 k\log(n(R+1))+C^2\log(n(C+1))=o(n),                            \tag{56}
\]
and the displayed catalog has only a polynomial factor in the input length.
\end{theorem}

\begin{proof}
The coordinate half-integer hyperplanes meeting the radius-$R$ ball in $U$
number at most $N$.  Their $k$-dimensional arrangement, including all tie
faces, has at most $2^k\sum_{j=0}^k\binom Nj$ rounding labels.

For an integer $z\in\cE_d$, put $x=P_Uz$, $u=\operatorname{round}(x)$, and
$w=z-u$.  If $w_i\ne0$, then $z_i$ is not the nearest integer to $x_i$, so
$|z_i-x_i|\ge1/2$.  Equation (54) implies
\[
 \operatorname{wt}(w)\le4C^2,
 \qquad |w_i|\le C+1/2.
\]
Thus $z\bmod2$ is obtained from one arrangement label by adding one of the
integer details counted by the second factor in (55).  This proves the cover
and (56).  Rational spectral projectors may be certified by enlarging $R,C$;
enumeration is output-sensitive, so an enormous $R$ is not treated as
polynomial in $\log R$.
\end{proof}

The simplex-cancellation family from Section~12 has a one-dimensional
minimum Gram eigenspace, $R=\Theta(\sqrt n)$, and $C<1$.  Hence (55) is
polynomial even though its shortest parity is the all-ones word and the
Hamming sieve is vacuous.  The rectangular family also has $k=1$ with
$R=O(1)$ and $C=2^{-\Omega(n)}$.  The same stable lifting construction thus
handles both known obstructions.

\subsection{Exact factor transports}

The size of the parity catalog is not the only exact complexity measure.
The coefficient energy itself may have a narrow factorization.  Let
$G=B^TB$ and, for the length guess $d\geq\lambda$, put
\[
 R_i:=\left\lfloor d\,\norm{\operatorname{row}_i(B^{-1})}_2\right\rfloor,
 \qquad D_i:=\{-R_i,\ldots,R_i\}.                              \tag{57}
\]
Join $i$ and $j$ when $G_{ij}\ne0$.  For an elimination order $\pi$, let
$K_i(\pi)$ contain the variable eliminated at step $i$ and its remaining
neighbors, including earlier fill edges, and define
\[
 W_\pi(B,d):=\max_i\sum_{j\in K_i(\pi)}\log_2(2R_j+1).          \tag{58}
\]

\begin{theorem}[Weighted Gram--GDL elimination]                 \label{thm:gdl}
For every explicitly supplied order $\pi$, exact search-$\SVP$ is solvable
in space
\[
 2^{W_\pi(B,d)}\poly(n,\operatorname{size}(B),\log(d+1))
\]
and at most $n$ times this quantity in time.  The algorithm returns a
minimizing coefficient vector by a backward-preimage pass.
\end{theorem}

\begin{proof}
If $v=Bz$ and $\norm v\le d$, then
\[
 |z_i|=|\ip{\operatorname{row}_i(B^{-1})}{v}|
 \le d\norm{\operatorname{row}_i(B^{-1})}_2,
\]
so every shortest coefficient vector lies in $\prod_iD_i$.  On this finite
product,
\[
 \norm{Bz}^2=\sum_iG_{ii}z_i^2+
              \sum_{i<j}2G_{ij}z_iz_j                         \tag{59}
\]
is a unary-and-pairwise factor graph with exactly the displayed Gram
interactions.  Exact min--sum elimination in order $\pi$ forms a table only
on the current bag $K_i(\pi)$, hence at most $2^{W_\pi(B,d)}$ entries.
Store an attaining value of the eliminated variable at every table entry;
the reverse pass recovers a global minimizer.  Exact rational arithmetic has
polynomial bit growth in these operations.

To exclude zero, run the elimination once for each $p\in[n]$ with the domain
$D_p\setminus\{0\}$ and take the best answer.  Every nonzero assignment
appears in at least one run, and a shortest assignment lies in the bounded
product.  This proves the claim.
\end{proof}

The same proof applies to auxiliary exact factorizations.  Suppose finite
variables $x$ have domains $E_x$, local energy factors and zero--infinity
constraint factors represent (59), and projection of the finite-objective
assignments onto $z$ is exactly $\prod_iD_i$.  A tree decomposition
$\mathcal D$ has weighted junction width
\[
 J(\mathcal D):=\max_{A\in\mathcal D}
                 \sum_{x\in A}\log_2|E_x|.                    \tag{60}
\]
Min--sum and backward preimage then cost $2^{J(\mathcal D)}\poly$ up to the
same $n$ zero-exclusion runs.  For example, after clearing denominators, a
row $(Bz)_k=\sum_iB_{ki}z_i$ can be represented by a binary tree of exact
partial sums and one unary square factor.  A sum over $S$ needs only the
integer interval
\[
 \left[-\sum_{i\in S}|B_{ki}|R_i,
        \sum_{i\in S}|B_{ki}|R_i\right].                       \tag{61}
\]
Thus a dense Gram graph need not imply a wide exact computation.  Conversely
the state intervals in (61) can be expensive; the certified quantity is the
weighted junction width, not unweighted treewidth.

\begin{corollary}[Exact-factor escape]                           \label{cor:gdl}
If an efficiently constructed order or auxiliary decomposition has width at
most $(1/2+o(1))n$, it gives an unconditional half-exponential algorithm on
that input.  In particular, width below
$(0.0264828-\varepsilon)n$ disposes of the input well inside the later
parity-query allowance.  No optimization over undisplayed bases or
decompositions is assumed.
\end{corollary}

There is also an exact dense-factor escape.  It is classical in the language
of lattices of Voronoi's first kind \cite{McKilliamGrant12}; the short proof
is included because its coarea form is the relevant transport principle.

\begin{theorem}[Obtuse-superbase minimum cut]                   \label{thm:mincut}
Let $b_1,\ldots,b_{n+1}$ be a superbase: $b_1,\ldots,b_n$ is a
lattice basis, $\sum_i b_i=0$, and
$\ip{b_i}{b_j}\le0$ for $i\ne j$.  Put
$w_{ij}=-\ip{b_i}{b_j}\ge0$.  A global minimum nontrivial cut of the graph
with weights $w_{ij}$ returns a globally shortest lattice vector.
\end{theorem}

\begin{proof}
Coefficients $c_i\in\Z$ are defined modulo addition of a common integer, and
the graph-Laplacian identity gives
\[
 \norm{\sum_i c_ib_i}^2
   =\sum_{i<j}w_{ij}(c_i-c_j)^2.                               \tag{62}
\]
For a nonzero vector normalize $\min_i c_i=0$ and put
$S_t=\{i:c_i\ge t\}$ for $1\le t\le\max_i c_i$.  Since an integer $a$
satisfies $a^2\ge|a|$,
\[
 \norm{\sum_i c_ib_i}^2
 \ge\sum_t\sum_{i<j}w_{ij}
       |\mathbf1_{i\in S_t}-\mathbf1_{j\in S_t}|.             \tag{63}
\]
Every summand in $t$ is a nontrivial cut capacity, so (63) is at least the
global minimum cut.  Conversely, the subset vector $\sum_{i\in S}b_i$ has
squared norm exactly the capacity of the cut $S$.  The two minima are equal,
and a minimum cut supplies the desired vector.
\end{proof}

Obtuseness of a displayed rational superbase is exactly checkable and a
global minimum cut is polynomial-time computable.  In particular the
simplex-cancellation family used above has an obtuse displayed superbase:
its exponential subset fan is submodular and is not part of the remaining
hard core.

\subsection{A direct gap escape}

Fix an oriented shortest vector $v$, write $a=v/\lambda$, and let
\[
 \gamma(v):=\min\{\norm w/\lambda:w\in L,\ w\notin\R v\}.       \tag{64}
\]
Positive multiples of $v$ cannot win the first-exit competition before $v$.
Define
\[
 a_\gamma:=
 \begin{cases}
  \sqrt{1-\gamma^2/4},&1\le\gamma\le\sqrt2,\\
  1/\gamma,&\gamma\ge\sqrt2.
 \end{cases}                                                    \tag{65}
\]

\begin{lemma}[Gap-dependent winner cap]                         \label{lem:gap-cap}
If $u\in S^{n-1}$ satisfies $\ip{u}{a}\ge a_\gamma$, the first facet met by
the ray $\R_{ge0}u$ is $F_v$.  Consequently the two winner caps for
$\pm v$ have total mass
\[
 2^{-(\delta_{\rm gap}(\gamma)+o(1))n},\qquad
 \delta_{\rm gap}(\gamma)
   =-\tfrac12\log_2(1-a_\gamma^2).                             \tag{66}
\]
\end{lemma}

\begin{proof}
Write a noncollinear competitor as $w=q\lambda b$, with $\norm b=1$ and
$q\ge\gamma$.  Minimality of $v-w$ implies $\ip{a}{b}\le q/2$.
Put $s=\ip{u}{a}$.  For $q\le\sqrt2$, maximizing $\ip{u}{b}$ under this constraint
shows that $\ip{u}{b}/q\le s$ once
$s\ge\sqrt{1-q^2/4}$.  For $q\ge\sqrt2$, the bound
$\ip{u}{b}\le1$ gives the sufficient condition $s\ge1/q$.  Both thresholds
decrease in $q$, yielding (65).  The fixed-height spherical-cap asymptotic is
$(1-a_\gamma^2)^{n/2}$ up to subexponential factors.
\end{proof}

At the query allowance in (34), this lemma closes the ledger whenever
\[
 \gamma>(1-2^{-2(0.0264828\ldots)})^{-1/2}
       =5.2670066\ldots.                                      \tag{67}
\]
The second branch in (65) is essential; continuing the first branch beyond
$\sqrt2$ would give an invalid constant.

\subsection{Causal transport through the Voronoi tiling}

Let $V$ be the origin Voronoi cell and let $Q$ be nearest-lattice decoding
away from its measure-zero ties.  For an oriented Voronoi-relevant vector
$w$, write
\[
 F_w:=V\cap\{x:\ip{x}{w}=\norm w^2/2\},\qquad
 A_w:=\mathcal H^{n-1}(F_w).                                  \tag{68}
\]
For a parity $\theta$, let $A_\theta$ sum these areas over the facets whose
coefficient vectors have parity $\theta$.  Draw $X$ uniformly from $V$ and
$U$ uniformly from $S^{n-1}$, independently, and set
\[
 C_\eta:=Q(X+\eta\lambda U)-Q(X).                              \tag{69}
\]

\begin{theorem}[Causal tiling kernel and transport conservation]
                                                                    \label{thm:causal}
For every nonzero parity $\theta$,
\[
 \lim_{\eta\to0+}
 \frac{\Prb[C_\eta\bmod2L=\theta]}{\eta}
 =d_n\frac{\lambda A_\theta}{\det L},\qquad
 d_n:=\frac{\Gamma(n/2)}{2\sqrt\pi\,\Gamma((n+1)/2)}
     =\Theta(n^{-1/2}).                                       \tag{70}
\]
Moreover, with
\[
 q_w:=\frac{\norm w A_w}{2n\det L},                            \tag{71}
\]
where both orientations are included,
\[
 \sum_wq_w=1,\qquad
 \frac1{2\det L}\sum_w A_w\frac{ww^T}{\norm w}=I_n,\qquad
 \sum_wq_w\frac{ww^T}{\norm w^2}=\frac1nI_n.                 \tag{72}
\]
\end{theorem}

\begin{proof}
Outside a tube of volume $O(\eta^2)$ around the codimension-two faces, a
segment of length $\eta\lambda$ crosses at most one facet.  At a regular
point of $F_w$, the starting points in $V$ which cross outward have normal
thickness
\[
 \eta\lambda\left(\ip{U}{w/\norm w}\right)_++O(\eta^2).
\]
Integrate over $F_w$, average over $U$, divide by $\operatorname{vol}V=\det
L$, and use $\E(U_1)_+=d_n$.  Summing the facets of parity $\theta$ proves
(70).

The pyramids $\operatorname{conv}(0,F_w)$ partition $V$ up to measure zero;
their volumes are $\norm wA_w/(2n)$.  This proves the first identity in
(72).  For the matrix identity, put $R(x)=x-Q(x)$.  Haar translation
invariance on $\R^n/L$ gives
\[
 \E[Q(X+h)-Q(X)]=h-\E[R(X+h)-R(X)]=h.                          \tag{73}
\]
Differentiate at $h=0$.  Crossing $F_w$ has one-sided rate
$A_w(\ip{h}{w/\norm w})_+/\det L$ and transports the label by $w$.
Pairing $w$ and $-w$ converts the positive parts into the middle identity of
(72).  Multiplying each summand by the definition (71) gives the last
identity; taking the trace also recovers $\sum_wq_w=1$.
\end{proof}

This realizes the principle that information lives in transport: the label
differences across actual first-arrival cells form an exact isotropic frame.
If $q_{\min}$ is the sum of (71) over the oriented shortest facets, (70)
gives total shortest-parity derivative at least
\[
 2nd_nq_{\min}=\Theta(\sqrt n)q_{\min}.                         \tag{74}
\]
One may draw $X$ uniformly in any fundamental parallelepiped and subtract
the nearest labels before and after the step; translation covariance gives
the same law.  This is a geometric representation, not an available $\CVP$
oracle.

The isotropic identity by itself cannot lower-bound $q_{\min}$.

\begin{proposition}[Isotropy alone has no atom bound]           \label{prop:isotropy-limit}
For every $n\ge2$ and every $0<\varepsilon<1/n$, there is a finite centrally
symmetric probability frame on $S^{n-1}$ with covariance $I_n/n$ but with
total weight exactly $\varepsilon$ on a designated antipodal pair.
\end{proposition}

\begin{proof}
Fix a unit vector $u$, an orthonormal basis $e_2,\ldots,e_n$ of $u^\perp$,
and put
\[
 \alpha=\frac{1/n-\varepsilon}{1-\varepsilon}.
\]
Give $\pm u$ total weight $\varepsilon$.  Distribute the remaining weight
uniformly over the centrally symmetric vectors
\[
 \pm\bigl(\sqrt\alpha\,u\pm\sqrt{1-\alpha}\,e_j\bigr),
 \qquad 2\le j\le n.
\]
The cross terms cancel.  The covariance in the $u$ direction is
$\varepsilon+(1-\varepsilon)\alpha=1/n$ and that on $u^\perp$ is
$(1-\varepsilon)(1-\alpha)/(n-1)=1/n$.
\end{proof}

Thus a proof of the remaining theorem must use the fact that (72) is the
cone-volume measure of a lattice Voronoi parallelohedron, or must combine it
with the exact escapes above; frame conservation alone cannot suffice.

There is a second obstruction to any proof that uses only local shortest
geometry.  Let $p_w$ be the spherical measure of rays exiting $V$ through
$F_w$, and let $p_{\min}$ sum these weights over the oriented shortest
facets.

\begin{proposition}[Kissing perturbation stress test]           \label{prop:kissing-stress}
Suppose a lattice $L$ has $K$ oriented shortest vectors.  For every
$\zeta>0$ there is an arbitrarily small linear perturbation $L'$ for which
one antipodal pair is uniquely shortest and
\[
 p_{\min}(L')+q_{\min}(L')\le\frac{4+\zeta}{K}.
\]
If $L$ is rational, the perturbation may be chosen rational.
\end{proposition}

\begin{proof}
Group the $K$ oriented shortest vectors into $K/2$ antipodal pairs.  Their
total direct-exit mass is at most one and their total cone-volume mass is at
most one.  Some pair $\{\pm v\}$ therefore has combined $p$- and $q$-weight
at most $4/K$.

Fix a shortest vector $v$, put $a=v/\lambda$, and apply
$T_\tau=I-\tau aa^T$ for sufficiently small $\tau>0$.  Among the old
shortest vectors, $\pm v$ shrink strictly more than every other pair, by the
strict case of Cauchy--Schwarz.  Lattice discreteness gives a positive gap to
the next length shell, so for small enough $\tau$ no longer vector overtakes
them.  Thus $\pm T_\tau v$ are uniquely shortest.

Every shortest facet contains the inball proved below, hence persists under
a sufficiently small perturbation.  Its solid angle, area, support distance,
and the Voronoi-cell volume vary continuously.  The two oriented weights of
the new shortest pair therefore converge to their two old weights as
$\tau\to0$, giving the stated bound.  When $L$ is rational,
$aa^T=vv^T/\lambda^2$ is rational, so rational $\tau$ preserves rationality.
\end{proof}

Consequently any exponential kissing family with exponent $\kappa$ would
rule out a universal cone-volume bound with $\delta<\kappa$, unless its small
perturbations are removed by one of the exact joint-core escapes.  This is
why code-derived high-kissing configurations are a necessary stress test for
the remaining theorem; a spherical-code bound alone does not decide whether
the corresponding lattice inputs survive those escapes.

\subsection{A log-theta collision kernel}

The causal crossing need not occur near the center $v/2$ of its facet.  This
prevents direct use of the raw theta Hessian: its two nearest terms retain a
tangential rank-one component.  The logarithmic Hessian cancels that
component exactly.  Put
\[
 Z_s(x):=\rho_s(L+x)=\sum_{w\in L}e^{-\pi\norm{x-w}^2/s^2}
\]
and define the posterior on lattice sites
\[
 \Prb_x[W=w]:=\frac{e^{-\pi\norm{x-w}^2/s^2}}{Z_s(x)}.          \tag{75}
\]

\begin{theorem}[Log-theta covariance and two-site stability]
                                                                    \label{thm:log-collision}
For every $x\in\R^n$,
\[
 \mathcal C_s(x):=
 \frac{s^4}{4\pi^2}\left(\nabla^2\log Z_s(x)+\frac{2\pi}{s^2}I_n\right)
 =\operatorname{Cov}_x(W).                                    \tag{76}
\]
Let $x\in F_v$ for a shortest vector $v$, set
$a=e^{-\pi\norm x^2/s^2}=e^{-\pi\norm{x-v}^2/s^2}$, and define
\[
 \epsilon_0(x):=\frac1{2a}
   \sum_{w\notin\{0,v\}}e^{-\pi\norm{x-w}^2/s^2},\qquad
 \epsilon_2(x):=\frac1{2a\lambda^2}
   \sum_{w\notin\{0,v\}}\norm w^2e^{-\pi\norm{x-w}^2/s^2}.      \tag{77}
\]
If $\epsilon_0(x),\epsilon_2(x)\le\epsilon\le1$, then
\[
 \norm{\mathcal C_s(x)-\tfrac14vv^T}_{\op}
 \le6\epsilon\lambda^2.                                       \tag{78}
\]
Consequently, when $\epsilon=o(1)$ the leading eigendirection of
$\mathcal C_s(x)$ is $o(1)$-close to $\R v$.
\end{theorem}

\begin{proof}
Differentiating the log partition function gives
\[
 \nabla\log Z_s(x)=-\frac{2\pi}{s^2}(x-\E_xW)
\]
and a second differentiation gives (76).

For (78), divide all lattice sites by $\lambda$ and put $u=v/\lambda$.
Relative to the total weight $2a$ of $\{0,v\}$, let $r_1$ and $r_2$ be the
first and second raw moments of the remaining sites.  Cauchy--Schwarz gives
$\norm{r_1}\le\sqrt{\epsilon_0\epsilon_2}$ and
$\norm{r_2}_{\op}\le\epsilon_2$.  If $m=\E_x(W/\lambda)$, then
\[
 \norm{m-u/2}\le\sqrt{\epsilon_0\epsilon_2}+\epsilon_0/2
 \le3\epsilon/2,
\]
while the normalized second raw moment differs from $uu^T/2$ by at most
$\epsilon_2+\epsilon_0/2$.  Subtracting the two outer products bounds the
covariance error by
\[
 \epsilon_2+\epsilon_0/2+3\epsilon/2+(3\epsilon/2)^2
 \le6\epsilon.
\]
Rescaling by $\lambda^2$ proves (78).  Standard eigenvector perturbation
then proves the last assertion.
\end{proof}

There is an exact line-scan interpretation.  Suppose
$x(t)=x_0+t\xi$ crosses $F_v$ at $t=t_0$.  Ignoring the remainder in (77),
the posterior probability of the site $v$ is
\[
 p(t)=\left(1+
 \exp\!\left[-\frac{2\pi\ip{\xi}{v}(t-t_0)}{s^2}\right]\right)^{-1},
                                                                    \tag{79}
\]
and
\[
 \mathcal C_s(x(t))=p(t)(1-p(t))vv^T.                           \tag{80}
\]
Thus its leading eigenvalue has a logistic ridge centered exactly at the
facet, of physical width $s^2/|\ip{\xi}{v}|$.  The proof above, with the two
distinguished weights no longer assumed equal, shows that (80) has
operator error at most $10\epsilon\lambda^2$ whenever the analogues of
$\epsilon_0,\epsilon_2$ relative to their combined weight are at most
$\epsilon$.  If the crossing angle and $\epsilon^{-1}$ are polynomial and
$s^2/\lambda^2=1/\poly(n)$, a polynomial grid resolves this ridge and its
eigendirection.

The quantity $\mathcal C_s$ is algebraically obtainable from
$Z_s,\nabla Z_s,\nabla^2Z_s$, all encoded by the common DGS Fourier panel.
However, this normalization converts the existing absolute empirical error
into relative error.  The analytic part of the remaining theorem is
therefore reduced to two explicit estimates on a robust shortest-facet core:
\[
 \epsilon_0+\epsilon_2\le n^{-\Omega(1)}
 \quad\hbox{and}\quad
 Z_s(x)\ \hbox{is above the empirical relative-error floor}.   \tag{81}
\]
Neither estimate follows from nearest-pair geometry alone.

\subsection{Robust shortest-facet cores}

\begin{lemma}[Shortest-facet robust core]                       \label{lem:facet-core}
Every shortest facet $F_v$ contains, in its affine hyperplane, the ball of
radius $\lambda/(2\sqrt3)$ centered at $v/2$.  For $0<\varepsilon<1$, the
homothetic core
\[
 F_v(\varepsilon):=(1-\varepsilon)F_v+\varepsilon(v/2)          \tag{82}
\]
has area $(1-\varepsilon)^{n-1}A_v$, is at tangential distance at least
$\varepsilon\lambda/(2\sqrt3)$ from every ridge, and every $x$ in the core
satisfies
\[
 \norm{x-w}^2-\norm x^2\ge\varepsilon\lambda^2/2               \tag{83}
\]
for all $w\in L\setminus\{0,v\}$.
\end{lemma}

\begin{proof}
On the hyperplane of $F_v$, the distance from $v/2$ to the constraint
induced by $w$ is
\[
 \frac{\norm w^2-\ip{v}{w}}{2\norm{P_{v^\perp}w}}.             \tag{84}
\]
Scale $\lambda=1$, set $A=\norm w^2$ and $b=\ip{v}{w}$.  Minimality and
Cauchy--Schwarz give $A\ge1$, $b\le A/2$, and $b^2\le A$.  On this domain
\[
 3(A-b)^2\ge A-b^2.                                            \tag{85}
\]
Indeed, for $1\le A\le4$ the difference is minimized at $b=A/2$,
where it equals $A(A-1)$; for $A\ge4$ it is minimized at $b=\sqrt A$.
Equations (84)--(85) give the inball.

Every defining slack of $F_v$ is affine and nonnegative on $F_v$.  At $v/2$
the slack belonging to $w\ne0,v$ is
$\norm w^2-\ip{v}{w}\ge\norm w^2/2\ge\lambda^2/2$.  The convex
combination (82) therefore gives (83), while applying it to the inball gives
the ridge distance.  Homothety scales $(n-1)$-area by
$(1-\varepsilon)^{n-1}$.
\end{proof}

Taking $\varepsilon=1/n$ retains a constant fraction of every shortest-facet
area and gives inverse-polynomial positional reach and energy margin.  Thus
robustification costs no exponential factor.  The inball also yields the
following unconditional, though generally weak, density escape.  If
$N_{\min}$ counts oriented shortest facets and $\kappa_m$ is the volume of
the unit $m$-ball, then
\[
 q_{\min}\ge
 \frac{N_{\min}\kappa_{n-1}\lambda^n}
      {2n(2\sqrt3)^{n-1}\det L}.                               \tag{86}
\]

\subsection{Tail suppression and the response-floor barrier}

The robust core gives a clean weighted-tail estimate, but its parameter
tradeoff is incompatible with the desired mass.  We state the obstruction
for Hhan's normalization (6).

\begin{lemma}[Homothetic weighted-tail bound]                  \label{lem:tail-bound}
Fix $t>0$ and a constant $0<\eta\le1$, put $s=1/\xi_t(d)$, and
let $x\in F_v(\eta)$.  The tail quantities in (77) satisfy
\[
 \epsilon_0(x)+\epsilon_2(x)
 \le 2^{(\frac12\log_2(t_0/(\eta t))+o(1))n}.                 \tag{87}
\]
Consequently the universal shell estimate certifies polynomial (indeed,
exponential) suppression for fixed parameters when $\eta t>t_0$.  Since the
available target scale has $t<1/4$, this forces $\eta>4t_0=0.9258\ldots$ and
retains at most
\[
 (1-4t_0+o(1))^{n-1}A_v
   =2^{-(3.7540+o(1))n}A_v.                                   \tag{88}
\]
\end{lemma}

\begin{proof}
For $w\in L$, write
\[
 D_w(z):=\norm{z-w}^2-\norm z^2=\norm w^2-2\ip zw.
\]
This is affine in $z$.  It is nonnegative on $F_v$ for every
$w\notin\{0,v\}$, while at the midpoint
\[
 D_w(v/2)=\norm w^2-\ip vw
 =\frac{\norm w^2+\norm{w-v}^2-\lambda^2}{2}
 \ge\frac{\norm w^2}{2}.
\]
The homothetic definition therefore gives
$D_w(x)\ge\eta\norm w^2/2$.  Dividing every non-colliding term by the two-site
weight $2a$ in (77) yields
\[
 \epsilon_0(x)+\epsilon_2(x)
 \le\frac12\sum_{w\ne0}
 \left(1+\frac{\norm w^2}{\lambda^2}\right)
 \exp\left(-\frac{\pi\eta\norm w^2}{2s^2}\right).
\]
Apply Corollary~3.2 of \cite{Hhan26} with Gaussian dilation
$a_\eta=\sqrt{2/\eta}$ and moment degree two.  Its exponent is
$\frac12\log_2(a_\eta^2t_0/(2t))=\frac12\log_2(t_0/(\eta t))$,
which proves (87).  The area formula in Lemma~\ref{lem:facet-core} and
$t<1/4$ prove (88).
\end{proof}

The preceding loss concerns one convenient core.  The next bound applies
to an arbitrary subset of a shortest facet and exposes the underlying
conflict.  A normalized theta response estimated from
$M=2^{(m+o(1))n}$ bounded Fourier samples has the usual absolute scale
$M^{-1/2}=2^{-(m/2+o(1))n}$.  In fact, for
$X\sim D_{L^*,1/s}$ the direct empirical Fourier summand obeys
\[
 \operatorname{Var}\bigl(\cos(2\pi\ip Xx)\bigr)
 =\frac{1+F_s(2x)}2-F_s(x)^2.
\]
Thus its variance is bounded below by a constant whenever $F_s(x)=o(1)$;
the $M^{-1/2}$ scale is present rather than merely an artifact of a uniform
upper bound.  We deliberately allow arbitrary polynomial slack around it.

\begin{theorem}[Empirical response-floor barrier]              \label{thm:response-barrier}
Fix $t>0$, let $M=2^{(m+o(1))n}$, and put $\alpha=m/2$ and
$s=1/\xi_t(d)$.  For a shortest facet $F_v$, let $E_v\subseteq F_v$ be any
set on which, for fixed $c,C>0$,
\[
 \epsilon_0(x)\le n^{-c},\qquad
 F_s(x)\ge n^{-C}2^{-\alpha n}.                               \tag{89}
\]
Writing $x=v/2+y$ with $y\perp v$, every $x\in E_v$ satisfies
\[
 \frac{\norm{y}^2}{\lambda^2}
 \le\frac14\left(\frac{\alpha}{t}-1\right)
       +O\left(\frac{\log n}{n}\right).                       \tag{90}
\]
If the leading term on the right lies in $(0,1/12)$, then
\[
 \frac{\mathcal H^{n-1}(E_v)}{A_v}
 \le
 \left(\sqrt{3(\alpha/t-1)+O(\log n/n)}\right)^{n-1}.          \tag{91}
\]
If $\alpha\le t$, the fraction is instead at most
$2^{-\frac12n\log_2 n+O(n\log\log n)}$ (and can be empty).
\end{theorem}

\begin{proof}
On $F_v$, orthogonality gives equal distinguished weights
\[
 a=\exp\left[-\pi\xi_t(d)^2
       \left(\frac{\lambda^2}{4}+\norm{y}^2\right)\right].
\]
The first condition in (89) implies $Z_s(x)=2a(1+\epsilon_0(x))\le3a$ for
large $n$.  Since $\rho_s(L)\ge1$, one has $F_s(x)=Z_s(x)/\rho_s(L)\le3a$.
Combining this with the response floor in (89), taking base-two logarithms,
and using $\xi_t(d)^2=4nt\ln2/(\pi d^2)$ gives
\[
 t\left(\frac{\lambda^2}{d^2}
       +\frac{4\norm{y}^2}{d^2}\right)
 \le\alpha+O(\log n/n).
\]
Because $\lambda\le d\le(1+1/n)\lambda$, this is (90).

Thus $E_v$ lies in an $(n-1)$-ball about $v/2$ of the radius in (90).
Lemma~\ref{lem:facet-core} lower-bounds $A_v$ by the volume of the
$(n-1)$-ball of radius $\lambda/(2\sqrt3)$.  Dividing the two ball volumes
proves (91).  When $\alpha\le t$, (90) gives radius
$O(\lambda\sqrt{\log n/n})$; the same volume comparison gives the stated
superexponential bound.
\end{proof}

Even granting the entire $2^{n/2+o(n)}$ sample panel to one query gives
$m=1/2$ and $\alpha=1/4$.  Tail suppression from Lemma~\ref{lem:tail-bound}
requires $t>t_0$ even at $\eta=1$.  Letting $t\downarrow t_0$ in (91) shows
that the accessible fraction of a shortest facet is at most
$2^{-(1.0289+o(1))n}$.  Conversely, retaining fraction
$2^{-(\delta+o(1))n}$ in (91) requires
\[
 m\ge2t_0\left(1+\frac{2^{-2\delta}}{3}\right)
   =0.61169\ldots\qquad(\delta=0.0264828),                     \tag{92}
\]
strictly above the half-exponential budget.  The actual multi-query ledger
is smaller still.  Therefore the two estimates in (81) cannot hold on
enough cone-volume mass for the proposed empirical log-theta line scan.

\subsection{The remaining joint-core theorem}

Let $p_{\min}$ denote the spherical measure of directions whose first
Voronoi exit has a shortest parity.  At allowance
$\delta<0.0264828\ldots$, fix an explicit enumerable catalogue of
certificates and call an input a \emph{joint core relative to that catalogue}
after every certified short-dual quotient and basis preprocessing step if:
all Hamming, spectral-tube, and affine parity covers have exponent exceeding
$\delta$; every efficiently displayed exact factor decomposition has cost
exceeding the half-exponential budget in Corollary~\ref{cor:gdl}; and no
certified obtuse-superbase representation applies.

This qualification is essential.  Without a fixed certificate catalogue,
``joint core'' is not an intrinsic class of lattices and the proposed
inequality is not a well-posed universal assertion.  Even after fixing the
catalogue, Proposition~\ref{prop:kissing-stress} shows that neither isotropy
nor local shortest-facet geometry can prove the desired exponent.

The surviving geometric statement is now the following, and is deliberately
recorded as a conjecture rather than a theorem.

\begin{conjecture}[Joint-core cone-volume bound]                 \label{conj:joint-core}
For some universal $\delta<0.0264828\ldots$, every joint core satisfies
\[
 p_{\min}+q_{\min}\ge2^{-(\delta+o(1))n}.                      \tag{93}
\]
\end{conjecture}

Theorem~\ref{thm:causal} and Lemma~\ref{lem:facet-core} settle the geometric
conversion from $q_{\min}$ to a robust causal crossing set, losing only a
polynomial factor.  Proposition~\ref{prop:isotropy-limit} shows why the
matrix conservation law does not prove (93) on its own.
Theorem~\ref{thm:log-collision} supplies the exact line observable and removes
all tangential-location dependence.  Theorem~\ref{thm:response-barrier},
however, proves that even (93) would not instantiate the empirical line scan:
the tail and response-floor estimates (81) cannot coexist on the required
mass under the stated sample budget.  The genuinely viable target is
therefore a lower bound on \emph{response-weighted} shortest-collision mass,
together with an estimator exploiting its variance, or a different nonlocal
observable that avoids division by the exponentially small local theta
response.  Direct-plus-cone mass alone is no longer a sufficient theorem.

\section{Conclusion}

Hhan's midpoint-Hessian theorem already supplies the exact recovery endpoint.
The Fourier samples used there contain a larger continuous object: a
traceless matrix field and its derivatives on the whole torus.  Bounded
narrowing importance sampling exposes this field at a rigorously controlled
smooth scale with exponent $0.4002$, log-trace-exponential smoothing makes
its gradient gap-free, and half-grid stationarity turns a successful
continuous endpoint into an exact target midpoint without branch transport.

These facts prove a black-box reduction from robust shortest-parity basin
mass to exact $\SVP$ in $2^{n/2+o(n)}$ time and space.  They do not prove that
the universal soft-Hessian flow supplies that mass.  On the contrary, the
anisotropic product family proves an essentially $2^{-n}$ upper bound for
the proposed Haar-seeded law and resolves the former Conjecture~13.1
negatively.

That counterexample does not obstruct every lattice-adapted seed.  Exact
short-dual equations remove irrelevant parity bits, the coefficient norm
gives a sparse Hamming law, and the spectral-tube theorem carries a
low-dimensional coarse response together with only a sparse lifting detail.
Weighted min--sum elimination gives a different exact escape whenever the
coefficient energy has a narrow displayed factorization.  An obtuse
superbase gives a polynomial minimum-cut algorithm, so the
simplex-cancellation family is solved for a stronger reason than its
spectral cover.

For inputs surviving those reductions, causal Voronoi transport supplies an
exact geometric kernel.  Its facet labels carry the crossing parities, its
cone-volume weights form an isotropic transport frame, and shortest facets
retain robust constant-area cores.  A sufficiently large noncollinear gap is
also solved by an explicit direct winner cap.  What remains is not the former
dense spectral core but a response-weighted version of the joint core.
Conjecture~\ref{conj:joint-core} remains a geometric question only after its
certificate catalogue is fixed; in any case, it is no longer sufficient for
the algorithm.  The response-floor barrier proves that the exact log-theta
covariance detector cannot combine polynomial tail suppression with the
required facet mass in $2^{n/2+o(n)}$ samples.  The lower constant in the
title is therefore conditional on a different observable or a variance-aware
response-weighted collision theorem.  The present paper proves the
surrounding reduction, exact escapes, and this obstruction, but not that
replacement clause.

\begin{thebibliography}{9}
\raggedright

\bibitem{ACKS25}
Divesh Aggarwal, Yanlin Chen, Rajendra Kumar, and Yixin Shen.
\newblock Improved classical and quantum algorithms for the shortest vector
problem via bounded distance decoding.
\newblock \emph{SIAM Journal on Computing}, 54(2):233--278, 2025.

\bibitem{ADRS15}
Divesh Aggarwal, Daniel Dadush, Oded Regev, and Noah
Stephens-Davidowitz.
\newblock Solving the shortest vector problem in $2^n$ time using discrete
Gaussian sampling.
\newblock In \emph{Proceedings of STOC 2015}, pages 733--742, 2015.

\bibitem{DRS14}
Daniel Dadush, Oded Regev, and Noah Stephens-Davidowitz.
\newblock On the closest vector problem with a distance guarantee.
\newblock In \emph{Proceedings of CCC 2014}, pages 98--109, 2014.

\bibitem{Hhan26}
Minki Hhan.
Solving the shortest vector problem in $2^{0.6039n}$ time via
mid-point Hessian.
arXiv:2608.02478v2, August 2026.

\bibitem{MR04}
Daniele Micciancio and Oded Regev.
\newblock Worst-case to average-case reductions based on Gaussian measures.
\newblock In \emph{Proceedings of FOCS 2004}, pages 372--381, 2004.

\bibitem{OpenAI26}
OpenAI.
\newblock Binary and spherical codes.
\newblock Chapter 2 of \emph{Ten Advances in Mathematics and Theoretical
Computer Science}, August 2026.
\newblock \url{https://cdn.openai.com/pdf/ten-proofs-oai.pdf}.

\bibitem{McKilliamGrant12}
Robert G. McKilliam and Alex Grant.
\newblock Finding short vectors in a lattice of Voronoi's first kind.
\newblock \emph{SIAM Journal on Discrete Mathematics},
28(3):1292--1302, 2014.
\newblock arXiv:1201.5154.

\bibitem{RSD17}
Oded Regev and Noah Stephens-Davidowitz.
\newblock An inequality for Gaussians on lattices.
\newblock \emph{SIAM Journal on Discrete Mathematics}, 31(2):749--757, 2017.

\end{thebibliography}

\end{document}
